\begin{abstract}
Resolver problemas visuais e l\^{o}gicos do benchmark ARC-AGI por meio da gera\c{c}\~{a}o autom\'{a}tica de c\'{o}digo em Python usando modelos de linguagem (LLMs) \'{e} um desafio central para a intelig\^{e}ncia artificial. Um problema recorrente nessa abordagem \'{e} o ajuste excessivo aos exemplos de treino (\textit{overfitting}): o modelo cria solu\c{c}\~{o}es que funcionam apenas para os exemplos mostrados, mas falham no teste. Neste trabalho, investigamos o m\'{e}todo de S\'{\i}ntese Indutiva Guiada por Contraexemplos (CEGIS) e avaliamos uma varia\c{c}\~{a}o com regras anti-ajuste (\textit{AntiCheat}), que pro\'{\i}be regras pontuais para pixels espec\'{\i}ficos e exige que o modelo explique em texto a l\'{o}gica geral antes de escrever o c\'{o}digo. Avaliamos 8 modelos de linguagem em 100 tarefas desafiadoras do ARC-AGI. Os resultados mostram que a inclus\~{a}o das regras anti-ajuste n\~{a}o alterou a acur\'{a}cia global de forma estatisticamente significante (acur\'{a}cia de 50,4\% no CEGIS contra 49,8\% no AntiCheat, com $p = 0,5515$). A an\'{a}lise qualitativa do c\'{o}digo gerado tamb\'{e}m indicou que ambos os m\'{e}todos produzem programas estruturalmente semelhantes e falham pelos mesmos motivos, principalmente pela incapacidade de encontrar a l\'{o}gica correta dentro do limite de tentativas. Por outro lado, a an\'{a}lise estrutural revelou que as tarefas seguem uma escala de dificuldade unidimensional quase perfeita de Guttman ($CR = 0,9738$) e que os c\'{o}digos com sobreajuste violam a Navalha de Occam, sendo 1,61 vezes maiores do que os programas que acertam de primeira ($p = 5,15 \times 10^{-9}$). Por fim, observamos que 86,0\% das corre\c{c}\~{o}es bem-sucedidas do CEGIS ocorrem at\'{e} a terceira itera\c{c}\~{a}o.
\end{abstract}

